\section{Methodology}
\label{sec:method}

\subsection{Hedonic rental-price model}

For listing $i$ of chip $c$ in month $t$,
\begin{equation}
\ln p_{ict} = \alpha_c + g(a_{ct}) + \delta_t + X_i'\beta + \varepsilon_{ict},
\label{eq:hedonic}
\end{equation}
where $p$ is the real posted price per GPU-hour, $a_{ct}$ the chip's age in years since its release
month, $g$ the age profile, $\alpha_c$ a chip (vintage) effect, $\delta_t$ a month effect and $X_i$
fixed effects for provider $\times$ region, contract type $\times$ term, GPUs per instance, price
basis (instance or accelerator only) and form factor. $g$ is a cubic spline in age with knots at
1.5, 3 and 5 years. The specifications were logged before any regression was run, together with
four hypotheses. H1, front-loaded decline: value falls fastest in a chip's first two years, closer
to geometric than to straight-line. H2, obsolescence shocks: a successor's release causes a
discrete drop in the predecessor's rent. H3, contract-type wedge: spot and on-demand prices fall
faster with age than reserved prices. H4, demand regimes: shortages mask age effects only
temporarily, so age profiles are stable across sub-periods once month effects are included.

\subsection{Identification}

Because $a_{ct} = t - v_c$, any linear trend in age can be moved into the chip and month effects,
so the linear part of $g$ is not identified in \eqref{eq:hedonic} \citep{hall1968technical}. Its
curvature is identified, but only through differences in calendar-time trends across vintages:
$a^2 = t^2 - 2 v t + v^2$, and $-2vt$ is a chip-specific trend. I report five objects.

\begin{enumerate}
\item \emph{Observed cohort paths.} Each chip's price path, by six-month age bin, with the controls
but without month effects. This needs no identification assumption: it is what happened.
\item \emph{Strategy 1, quality-adjusted vintage.} Replace $\alpha_c$ with observed specifications:
(a) $\ln$ BF16 dense TFLOP/s alone, and (b) adding $\ln$ memory bandwidth and $\ln$ memory. The
linear age term is then identified from the cross-chip relation between residual price and release
date.
\item \emph{Strategy 2, frontier-relative obsolescence.} Keep $\alpha_c$, drop $\delta_t$ and add the
chip's TFLOP/s relative to the best chip released by month $t$, after \citet{geske2007computers}.
The assumption is that market-wide time effects work only through the frontier.
\item \emph{Strategy 3, successor event study.} Event-time dummies around the release of each
chip's successor (V100 at A100, A100 at H100, H100 at B200), with chip and month effects.
\item \emph{Strategy 4, shape only.} Keep both sets of effects, identify only the curvature, and
report $g_d(a) = d\,a + \tilde c(a)$ over a grid of assumed average rates $d$ from 0 to $-0.40$
log points a year, where $\tilde c$ is the curvature detrended over ages 0--6.
\end{enumerate}

The summary measures are $g(1)-g(0)$, $g(2)-g(0)$, $g(4)-g(0)$, the average annual log change over
years 0--2 and 2--4, and their difference, the \emph{early minus late rate}. H1 predicts a negative
early minus late rate. Strategy 4 reports only that
difference, the one measure that does not depend on the assumed linear term.

\subsection{Inference}

Age has one value per chip and month, so inference is at the chip level, with seven clusters. The
method specified in advance is the unrestricted wild cluster bootstrap with Webb's six-point weights and
9,999 draws \citep{cameron2008bootstrap, webb2023reworking}. Cluster-robust standard errors by
provider and two-way by provider and month are shown for comparison. After the adversarial review
(Appendix~\ref{app:review}) I add the restricted bootstrap (WCR), which imposes the null when generating draws and is
recommended when clusters are few and unequal in size \citep{mackinnon2017wild,
mackinnon2023cluster}, a test-inversion interval, and a Monte Carlo of interval coverage on this
panel's design.

\subsection{From rents to asset values}

The value of a chip at age $a$ is the present value of its remaining net rent,
\begin{equation}
V(a) = \sum_{m = 0}^{12(H - a) - 1} \frac{730 \, \max\{u\,p(a + m/12) - o,\ 0\} \, S(m)}{(1+r)^{(m+1)/12}},
\label{eq:value}
\end{equation}
summed over remaining months $m$ at 730 hours a month, with utilization $u$, operating cost per GPU-hour $o$ (system
power $\times$ PUE $\times$ electricity price, plus maintenance at 4\% of capital cost a year),
survival $S$ at a 2\% annual failure rate, discount rate $r$ and horizon $H$ of nine years, the
oldest age observed. The rent path is $p(a) = p_{\text{now}} \exp\{g(a) - g(a_{\text{now}})\}$:
each chip's current neocloud median on-demand rent moved along a profile $g$ from the realized
cohort paths or one of the strategies. I report $V(0)$ relative to capital cost, the share of value
lost by each year, the economic life $T^*$ (the first age at which $u p \le o$), and a
straight-line-equivalent life $L = 2\int_0^H V(a)/V(0)\,da$, the straight-line life with the same
average book value. The base case is $r = 10\%$, $u = 70\%$ and a PUE of 1.3; the grid is
$r \in \{8, 10, 12\}\%$, $u \in \{50, 70, 90\}\%$ and PUE $\in \{1.2, 1.4\}$.

\subsection{Earnings restatement}

Section~\ref{sec:accounting} enters each year's estimated GPU capital expenditure as a purchase
cohort and depreciates it under the firm's disclosed straight-line life and under the
rental-implied schedules from \eqref{eq:value}, with the GPU share of capital expenditure varied
over a grid because no firm discloses it.
